\documentclass[11pt]{article}
\usepackage[margin=2.5cm]{geometry}
\usepackage{booktabs}
\usepackage{tabularx}
\usepackage[hidelinks]{hyperref}
\usepackage{times}

\title{New Embeddings That Fixed the Wrong Problem:\\ Experiments 08, 09 and 10 Rerun}
\author{}
\date{August 15, 2026 \quad (extended September 7, 2026)}

\begin{document}
\maketitle

\section*{The problem we started from}

Report 03 ended with one measurement: all the patient embeddings looked almost identical to
each other. Two random patients had a similarity of $.979$ out of $1$, so the ``20 nearest
neighbours'' in the graph were chosen by differences in the fourth decimal place.

The lab proposed a set of new encoders to fix this. This report describes them and what
happened when experiments 08, 09 and 10 were run again on top of them.

\section*{The new embeddings}

The old encoder was trained on one thing only: predict the label. Nothing stopped it from
throwing away every direction in the embedding except the one that helps with that prediction,
and that is what it did. The new recipes add a second term to the training objective that
pushes back: each of the 64 numbers in an embedding must keep some spread, and different
numbers must not repeat each other.

Five versions were trained. Each is scored two ways: how similar two random patients look
(lower is better spread), and how well the embeddings predict the label on their own, before
any graph is involved.

\begin{table}[h]
\centering
\begin{tabularx}{\textwidth}{lXcc}
\toprule
Version & What it does & Similarity & AUROC \\
\midrule
A & The old recipe, reproduced & $.981$ & $.798$ \\
B & A bigger encoder, same objective & $.890$ & $.769$ \\
C & Old encoder $+$ the anti-collapse term & $\mathbf{.633}$ & $\mathbf{.788}$ \\
D & Bigger encoder $+$ the anti-collapse term & $.356$ & $.742$ \\
E & Uses the patient's real visit sequence $+$ the term & $.338$ & $.710$ \\
\bottomrule
\end{tabularx}
\caption{The five encoders, one run each. Similarity is between two random patients.}
\end{table}

There is a clear trade: the more the embeddings are forced to spread out, the worse they
predict the label. Version C is the compromise --- it removes most of the crowding while
giving up only $.010$ of prediction quality --- so C is the one taken forward.

The fix works on its own terms. Random patients moved from $.979$ to $.633$, and the
similarities among the neighbour candidates, which is what actually decides the graph, went
from a range of $.009$ to a range of $.080$. The neighbours are now chosen by real
differences instead of rounding noise.

\section*{Experiment 08 again}

\begin{table}[h]
\centering
\begin{tabularx}{\textwidth}{Xcc}
\toprule
Model & Old embeddings & Version C \\
\midrule
rgat & $\mathbf{.800}$ & $.773$ \\
rgcn & $.784$ & $.768$ \\
\bottomrule
\end{tabularx}
\caption{Test AUC on the color task. Best setting chosen on validation, one run each.}
\end{table}

\paragraph{Everything got worse.} Both models lost between $.016$ and $.027$. Repeated runs of
the same setting differ by about $.004$--$.009$, so these drops are three to seven times
larger than the normal noise.

\paragraph{It is not bad luck in the tuning.} The whole set of tried settings moved down, not
just the best one: the middle setting went from $.783$ to $.770$. Validation scores dropped
by the same amount.

\paragraph{It is missing information, not overfitting.} The score on the \emph{training} data
fell as well, from $.866$ to $.803$. The model is not memorising more and generalising less;
there is simply less to work with.

\section*{Experiment 09 again}

\begin{table}[h]
\centering
\begin{tabularx}{\textwidth}{Xcc}
\toprule
Readout & Old embeddings & Version C \\
\midrule
\texttt{supernode} & $.795 \pm .009$ & $\mathbf{.780 \pm .008}$ \\
\texttt{pma} & $.795 \pm .004$ & $.776 \pm .007$ \\
\texttt{gated\_attn} & $.796 \pm .005$ & $.773 \pm .004$ \\
\texttt{node\_score\_mean} & $.796 \pm .002$ & $.772 \pm .004$ \\
\texttt{mean} (the old method) & $\mathbf{.798 \pm .004}$ & $.769 \pm .010$ \\
\texttt{multistat} & $.792 \pm .007$ & $.768 \pm .010$ \\
\bottomrule
\end{tabularx}
\caption{Test AUC on the color task, 5 runs each.}
\end{table}

\paragraph{The order changed.} On the old embeddings nothing beat the plain average and the
six numbers were indistinguishable. On the new ones the plain average drops to fifth place,
and \texttt{supernode} --- an extra node per patient group that gathers the information --- is
now $.011$ ahead of it. On the per-visit task the same gap is $.022$.

\paragraph{But this is a shuffle inside a lower range.} The best readout on the new embeddings
$(.780)$ is still below the \emph{worst} readout on the old ones $(.792)$.

\section*{Experiment 10 again}

\begin{table}[h]
\centering
\begin{tabularx}{\textwidth}{Xcc}
\toprule
Graph & Old embeddings & Version C \\
\midrule
Flexible & $\mathbf{.800 \pm .003}$ & $\mathbf{.780 \pm .005}$ \\
20 neighbours (the old graph) & $.796 \pm .002$ & $.767 \pm .004$ \\
\bottomrule
\end{tabularx}
\caption{Test AUC on the color task, rgat, 5 runs each.}
\end{table}

\paragraph{The flexible graph finally does something.} On the old embeddings it was $.004$
ahead of the fixed graph, small enough to be explained by the extra tuning it received. On the
new embeddings the gap is $.013$, about three times the run-to-run noise.

\paragraph{And it now prefers the opposite settings.} On the old embeddings the winner was a
dense graph (63 connections per patient) built on the untouched similarities. On the new ones
the winner is a sparse graph (12 connections) built on whitened similarities, and whitening is
now the best of the three ways of measuring similarity ($.794$, against $.780$ untouched and
$.770$ centred). The choice of similarity measure used to make no difference at all.

\paragraph{The gain does not close the gap.} The best result on the new embeddings is
$.780 \pm .005$, against $.800 \pm .003$ on the old ones.

\section*{What happens with no embeddings at all}

Experiments 08, 09 and 10 were run once more with the embeddings taken out of the model's
input. The graph was still built from them, so the same visits were still neighbours. Only the
vector each visit starts from changed. Four starting vectors were compared. This ran on the old
embeddings, so the number to compare with is $.800$.

\begin{table}[h]
\centering
\begin{tabularx}{\textwidth}{Xlcc}
\toprule
What each visit starts from & Visits start & Test AUC & On training data \\
\midrule
Its embedding: the 64 numbers the encoder
made from its text (every result above) & different, by text & $\mathbf{.800}$ & $.851$ \\
A random vector, drawn once before training
and never changed & different, by noise & $.702$ & $.831$ \\
A vector of its own that the model is free
to change while it trains & different, by noise & $.701$ & $1.000$ \\
One single vector, shared by every visit
in the dataset & identical & $.499$ & $.504$ \\
\bottomrule
\end{tabularx}
\caption{Experiment 08 rerun with the embeddings removed from the input. All four rows use the
same graph, so in all four the same visits are neighbours; only the starting vectors differ.
Test AUC on the color task, rgat, best setting chosen on validation, one run each. Chance is
$.500$.}
\end{table}

\paragraph{Without the embeddings the score falls from $.800$ to $.70$.} Chance is $.500$.

\paragraph{A trained vector and a random one score the same} --- $.701$ and $.702$. The trained
one scores a perfect $1.000$ on the training data, so it is memorising the visits it has seen.

\paragraph{One shared vector scores chance.} Every layer of the network averages over
neighbours, so visits that start the same stay the same.

\begin{table}[h]
\centering
\begin{tabularx}{\textwidth}{Xcc}
\toprule
 & With embeddings & Without \\
\midrule
Exp.\ 09: best of the six readouts & $.798 \pm .004$ (\texttt{mean}) & $.679 \pm .029$ (\texttt{pma}) \\
Exp.\ 09: worst of the six & $.792 \pm .007$ (\texttt{multistat}) & $.622 \pm .017$ (\texttt{supernode}) \\
\addlinespace
Exp.\ 10: flexible graph & $.800 \pm .003$ & $.705 \pm .020$ \\
Exp.\ 10: fixed graph, 20 neighbours each & $.796 \pm .002$ & $.670 \pm .029$ \\
\bottomrule
\end{tabularx}
\caption{The left column repeats tables 3 and 4; the right column is the same experiment with
the embeddings taken out of the input, using the trained vector from row 3 of table 5. The two
top rows are the best and the worst of the six readouts of experiment 09 --- which readout that
is differs between the columns, so each one is named. The two bottom rows are the two graphs of
experiment 10. What the text is about is the size of the gap between the rows of a pair, not the
rows themselves. Test AUC on the color task, rgat, 5 runs each.}
\end{table}

\paragraph{The choice of readout starts to matter.} The six readouts spread over $.057$ without
the embeddings, against $.006$ with them.

\paragraph{So does the choice of graph.} The flexible graph beats the fixed one by $.035$
without the embeddings, against $.004$ with them.

\paragraph{The winning graph is the same kind as before} --- dense, $55$ connections per
patient, built on the untouched similarities.

\section*{What this adds up to}

\paragraph{The crowding was not the bottleneck.} It was a real defect and it has been fixed,
but fixing it made every downstream number worse. The one direction the old encoder had
concentrated on was carrying the signal; spreading the embeddings out spent it.

\paragraph{The graph and the readout do respond now.} Both parts of the pipeline that looked
dead in report 03 came back to life on the new embeddings: the flexible graph, the choice of
similarity measure, and the choice of readout all matter where none of them did before. This
says those components were never broken --- they had nothing to work with.

\paragraph{The two effects do not cancel out.} The parts we can tune gained about $.013$; the
encoder change cost about $.030$. The best model overall is still the old one, at
$.800 \pm .003$.

\end{document}
